\documentclass[tikz,border=3pt]{standalone}
\usetikzlibrary{arrows.meta}
\begin{document}
\begin{tikzpicture}[>=Latex, font=\small]
  % ejes de fondo, tenues
  \draw[gray!50] (-1,0) -- (4.2,0);
  \draw[gray!50] (0,-0.5) -- (0,3.2);

  % W = span{(1,1)}
  \draw[thick] (-0.5,-0.5) -- (3.6,3.6) node[above right] {$W$};

  \coordinate (v) at (1.1,2.6);
  \coordinate (Tv) at (1.85,1.85); % pie de la perpendicular sobre W
  \coordinate (Uv) at (1.1,1.1);   % otra proyección sobre W

  \fill (v) circle (1pt) node[above] {$v$};
  \fill (Tv) circle (1pt) node[right] {$T(v)$};
  \fill (Uv) circle (1pt) node[below right] {$U(v)$};

  \draw[dashed] (v) -- (Tv);
  \draw[dashed] (v) -- (Uv);

  % marca de ángulo recto en T(v)
  \draw (Tv) ++(135:0.18) -- ++(-45:0.18) -- ++(-135:0.18);
\end{tikzpicture}
\end{document}
